\documentclass[10pt,a4paper]{amsart}
\usepackage{amsmath,amssymb,amsthm,graphicx}
\usepackage{verbatim}
\newenvironment{reference}{\comment}{\endcomment}
%\newenvironment{reference}{}{}
\newenvironment{slogan}{\comment}{\endcomment}
\newenvironment{history}{\comment}{\endcomment}
% For commutative diagrams we use Xy-pic
\usepackage[all]{xy}

% We use 2cell for 2-commutative diagrams.
\xyoption{2cell}
% Theorem environments.
%
\theoremstyle{plain}
\newtheorem{theorem}[subsection]{Theorem}
\newtheorem{proposition}[subsection]{Proposition}
\newtheorem{lemma}[subsection]{Lemma}

\theoremstyle{definition}
\newtheorem{definition}[subsection]{Definition}
\newtheorem{example}[subsection]{Example}
\newtheorem{exercise}[subsection]{Exercise}
\newtheorem{situation}[subsection]{Situation}

\theoremstyle{remark}
\newtheorem{remark}[subsection]{Remark}
\newtheorem{remarks}[subsection]{Remarks}

\numberwithin{equation}{subsection}

% Macros
%
\def\lim{\mathop{\mathrm{lim}}\nolimits}
\def\colim{\mathop{\mathrm{colim}}\nolimits}
\def\Spec{\mathop{\mathrm{Spec}}}
\def\Hom{\mathop{\mathrm{Hom}}\nolimits}
\def\Ext{\mathop{\mathrm{Ext}}\nolimits}
\def\SheafHom{\mathop{\mathcal{H}\!\mathit{om}}\nolimits}
\def\SheafExt{\mathop{\mathcal{E}\!\mathit{xt}}\nolimits}
\def\Sch{\mathit{Sch}}
\def\Mor{\mathop{\mathrm{Mor}}\nolimits}
\def\Ob{\mathop{\mathrm{Ob}}\nolimits}
\def\Sh{\mathop{\mathit{Sh}}\nolimits}
\def\NL{\mathop{N\!L}\nolimits}
\def\CH{\mathop{\mathrm{CH}}\nolimits}
\def\proetale{{pro\text{-}\acute{e}tale}}
\def\etale{{\acute{e}tale}}
\def\QCoh{\mathit{QCoh}}
\def\Ker{\mathop{\mathrm{Ker}}}
\def\Im{\mathop{\mathrm{Im}}}
\def\Coker{\mathop{\mathrm{Coker}}}
\def\Coim{\mathop{\mathrm{Coim}}}

% Boxtimes
%
\DeclareMathSymbol{\boxtimes}{\mathbin}{AMSa}{"02}

%
% Macros for moduli stacks/spaces
%
\def\QCohstack{\mathcal{QC}\!\mathit{oh}}
\def\Cohstack{\mathcal{C}\!\mathit{oh}}
\def\Spacesstack{\mathcal{S}\!\mathit{paces}}
\def\Quotfunctor{\mathrm{Quot}}
\def\Hilbfunctor{\mathrm{Hilb}}
\def\Curvesstack{\mathcal{C}\!\mathit{urves}}
\def\Polarizedstack{\mathcal{P}\!\mathit{olarized}}
\def\Complexesstack{\mathcal{C}\!\mathit{omplexes}}
% \Pic is the operator that assigns to X its picard group, usage \Pic(X)
% \Picardstack_{X/B} denotes the Picard stack of X over B
% \Picardfunctor_{X/B} denotes the Picard functor of X over B
\def\Pic{\mathop{\mathrm{Pic}}\nolimits}
\def\Picardstack{\mathcal{P}\!\mathit{ic}}
\def\Picardfunctor{\mathrm{Pic}}
\def\Deformationcategory{\mathcal{D}\!\mathit{ef}}
\usepackage{cleveref}
\crefname{originalreference}{original reference}{original references}
\Crefname{originalreference}{Original reference}{Original references}
\makeatletter
\newlabel{algebra-lemma-pullback-module}{{algebra-lemma-pullback-module}{source}{}{sourcepointer-1}{}}
\newlabel{algebra-lemma-pullback-module@cref}{{[originalreference][0][]algebra-lemma-pullback-module}{[1][0][]source}{}{}{}}
\makeatother
\crefname{originallemma}{original lemma}{original lemmas}
\Crefname{originallemma}{Original lemma}{Original lemmas}
\makeatletter
\newlabel{lemma-finite}{{lemma-finite}{source}{}{sourcepointer-2}{}}
\newlabel{lemma-finite@cref}{{[originallemma][0][]lemma-finite}{[1][0][]source}{}{}{}}
\makeatother
\crefname{originallemma}{original lemma}{original lemmas}
\Crefname{originallemma}{Original lemma}{Original lemmas}
\makeatletter
\newlabel{lemma-torsion-submodule-sum-injective-hulls}{{lemma-torsion-submodule-sum-injective-hulls}{source}{}{sourcepointer-3}{}}
\newlabel{lemma-torsion-submodule-sum-injective-hulls@cref}{{[originallemma][0][]lemma-torsion-submodule-sum-injective-hulls}{[1][0][]source}{}{}{}}
\makeatother
\crefname{originallemma}{original lemma}{original lemmas}
\Crefname{originallemma}{Original lemma}{Original lemmas}
\makeatletter
\newlabel{lemma-union-artinian}{{lemma-union-artinian}{source}{}{sourcepointer-4}{}}
\newlabel{lemma-union-artinian@cref}{{[originallemma][0][]lemma-union-artinian}{[1][0][]source}{}{}{}}
\makeatother
\title{The injective hull of the residue field base changes along a flat local homomorphism}
\author{The Stacks Project Authors}
\date{stacks-project:dualizing.tex:local-snapshot-20261012; GNU FDL 1.2 or later}
\begin{document}
\maketitle
\noindent\textit{Editorial note.} All delivered mathematical text is the author's own, copied from the pinned Stacks Project chapter source sources/stacks-dualizing/pinned-original-source.tex; no statement or proof step was rewritten, reordered or removed.\par\medskip
\noindent\textit{Editorial note.} Preamble provenance: the Stacks Project keeps one single shared preamble for every chapter; this card takes it directly from the staged sources/stacks-dualizing/preamble.tex as three verbatim spans (lines 8-12, 14-18 and 154-217), so no derived file was written and nothing was added or redefined. The wrapper class and the base packages are supplied by the assembler.\par\medskip
\noindent\textit{Editorial note.} Bibliography provenance: neither the counted statement nor its proof contains a citation (this lane rescans the whole span before writing the card), so the card carries no bibliography and no entry was reconstructed.\par\medskip
\noindent\textit{Editorial note.} Selection provenance: this card was opened under the file-based control contract (CONTROL.md), with the mechanical pre-card gate run on its statement and proof interval before the card was written.\par\medskip
\noindent\textit{Allowed standard prerequisites (editorial).}\begin{itemize}
\item The description of the hull as the union of the submodules killed by the powers of the maximal ideal. Invoked at the start of the counted proof (source line 5521). Reduces the statement to the truncated pieces over the Artinian quotients.
\item The equality of the lengths of the hull and of the ring over an Artinian local ring. Invoked inside the counted proof (source lines 5528-5531). Compares the length of the base changed hull with that of the hull of the target residue field.
\item The criterion for an injective module map between hulls to be an isomorphism. Invoked inside the counted proof (source lines 5540-5542). Produces the map and upgrades it to an isomorphism once the lengths agree.
\item The equality of lengths under the flat pullback of finite length modules. Invoked inside the counted proof (source lines 5525-5527). Makes the length of the base changed module computable over the target.
\end{itemize}\medskip
\ifdefined\hypertarget\hypertarget{sourcepointer-1}{}\fi
\noindent\textit{Source reference.} Original reference algebra-lemma-pullback-module: sources/stacks-dualizing/pinned-original-source.tex, source lines 5527--5527. Navigation pointer only: the counted text refers to the Stacks item algebra-lemma-pullback-module, which belongs to another chapter of the project and is not part of this card. The pointer carries the target's own label as its number because this chapter's staged snapshot does not contain the other chapter's tag table; its content is not reproduced here.\par\medskip
\ifdefined\hypertarget\hypertarget{sourcepointer-2}{}\fi
\noindent\textit{Source reference.} Original lemma lemma-finite: sources/stacks-dualizing/pinned-original-source.tex, source lines 912--912. Navigation pointer only: the counted text refers to the Stacks item lemma-finite, which is an existing labeled result of this same chapter and is not part of this card. The pointer carries the target's own label as its number and its definition line as its span; its content is not reproduced here.\par\medskip
\ifdefined\hypertarget\hypertarget{sourcepointer-3}{}\fi
\noindent\textit{Source reference.} Original lemma lemma-torsion-submodule-sum-injective-hulls: sources/stacks-dualizing/pinned-original-source.tex, source lines 1023--1023. Navigation pointer only: the counted text refers to the Stacks item lemma-torsion-submodule-sum-injective-hulls, which is an existing labeled result of this same chapter and is not part of this card. The pointer carries the target's own label as its number and its definition line as its span; its content is not reproduced here.\par\medskip
\ifdefined\hypertarget\hypertarget{sourcepointer-4}{}\fi
\noindent\textit{Source reference.} Original lemma lemma-union-artinian: sources/stacks-dualizing/pinned-original-source.tex, source lines 1042--1042. Navigation pointer only: the counted text refers to the Stacks item lemma-union-artinian, which is an existing labeled result of this same chapter and is not part of this card. The pointer carries the target's own label as its number and its definition line as its span; its content is not reproduced here.\par\medskip
\begin{lemma}
\label{lemma-injective-hull-goes-up}
Let $(A, \mathfrak m, \kappa) \to (B, \mathfrak n, l)$
be a flat local homorphism of Noetherian rings such that
$\mathfrak n = \mathfrak m B$. If $E$ is the injective
hull of $\kappa$, then $E \otimes_A B$ is the injective
hull of $l$.
\end{lemma}

\begin{proof}
Write $E = \bigcup E_n$ as in Lemma \ref{lemma-union-artinian}.
It suffices to show that $E_n \otimes_{A/\mathfrak m^n} B/\mathfrak n^n$
is the injective hull of $l$ over $B/\mathfrak n$.
This reduces us to the case where $A$ and $B$ are Artinian local.
Observe that $\text{length}_A(A) = \text{length}_B(B)$ and
$\text{length}_A(E) = \text{length}_B(E \otimes_A B)$
by Algebra, Lemma \ref{algebra-lemma-pullback-module}.
By Lemma \ref{lemma-finite} we have
$\text{length}_A(E) = \text{length}_A(A)$ and
$\text{length}_B(E') = \text{length}_B(B)$
where $E'$ is the injective hull of $l$ over $B$.
We conclude $\text{length}_B(E') = \text{length}_B(E \otimes_A B)$.
Observe that
$$
\dim_l((E \otimes_A B)[\mathfrak n]) =
\dim_l(E[\mathfrak m] \otimes_A B) =
\dim_\kappa(E[\mathfrak m]) = 1
$$
where we have used flatness of $A \to B$ and $\mathfrak n = \mathfrak mB$.
Thus there is an injective $B$-module map $E \otimes_A B \to E'$
by Lemma \ref{lemma-torsion-submodule-sum-injective-hulls}.
By equality of lengths shown above this is an isomorphism.
\end{proof}

\end{document}
